\section{Measuring and Predicting Stability Regimes}
\label{sec:method}

\paragraph{Setup.}
A policy $\pi$ maps an observation $o_t$ to an action chunk
$(a_t, \dots, a_{t+H-1})$ of nominal length $H$, here $H{=}16$. The executor
plays the first $k \le H$ actions open-loop and then replans. We ask how $k$
should depend on the state.

\paragraph{Premise: what an open-loop chunk costs.}
The price of executing $k$ actions blind can be written down before it is
measured. Take the standard incremental-stability model of the plant: the
dynamics are $(C,\rho)$-exponentially incrementally input-to-state stable if any
two trajectories driven by input sequences $u$ and $u'$ satisfy
\begin{equation}
\lVert x_j - x'_j \rVert \;\le\; C\rho^{\,j-1}\lVert x_1 - x'_1 \rVert
\;+\; C\sum_{i=1}^{j-1} \rho^{\,j-1-i}\,\lVert u_i - u'_i \rVert ,
\label{eq:eiss}
\end{equation}
with contraction rate $\rho$ and constant $C \ge 1$
\citep{simchowitz2025pitfalls,zhang2025chunking}. The two things a chunked
executor does are the two special cases of \eqref{eq:eiss}.

\emph{Perturb one input.} Start both branches at the same state and perturb only
$u_1$, by $\sigma_u$. Then \eqref{eq:eiss} collapses to a single term,
$\log \lVert x_j - x'_j \rVert \le \log C\sigma_u + (j-1)\log\rho$, a straight
line in $j$ of slope $\log\rho$. The labeling procedure below perturbs exactly
the first action and fits exactly that slope, so the $\lambda$ it returns is a
state-local empirical estimate of $\log\rho$ in units of inverse control steps,
and the linear fit is the model rather than a convenience.

\emph{Execute a chunk.} A policy whose per-step action error is $\varepsilon$
commits that error at every step of the chunk, since nothing corrects it until
the next replan. Putting $\lVert u_i - u'_i \rVert = \varepsilon$ in
\eqref{eq:eiss} and writing $\rho = e^{\lambda}$, the deviation at the end of a
chunk of length $k$ obeys
\begin{equation}
\lVert x_k - x'_k \rVert \;\le\; C\varepsilon\,
\frac{e^{\lambda k}-1}{e^{\lambda}-1}
\;=\;
\begin{cases}
C\varepsilon\,/\,(1-e^{\lambda}) & \lambda < 0,\\[2pt]
C\varepsilon\,k & \lambda = 0,\\[2pt]
C\varepsilon\,e^{\lambda k}/(e^{\lambda}-1) & \lambda > 0.
\end{cases}
\label{eq:chunk}
\end{equation}
Chunk length is free where the segment is open-loop stable, costs linearly at
the noise floor, and costs exponentially above it. The scale is not academic:
the two states in Figure~\ref{fig:labeling}(b) carry $\lambda{=}.006$ and
$\lambda{=}.121$, so a single 16-step chunk multiplies a given action error by
$1.1$ at the first and by $6.9$ at the second.

\emph{Over an episode.} Run $T/k$ chunks and let a replan contract the
accumulated deviation by $\gamma<1$ where the policy is closed-loop stable. The
deviation at chunk boundaries then follows a linear recursion of gain
$\gamma e^{\lambda k}$, which stays bounded over the episode exactly when
$k < \log(1/\gamma)/\lambda$. Open-loop instability sets a ceiling on chunk
length, and closed-loop stability is what raises that ceiling.

Two known results bracket this picture. With no assumption on the dynamics at
all, the classical imitation bound is
$J(\hat\pi) \le J(\pi^\star) + T^2\varepsilon$ for a policy with per-step error
$\varepsilon$, and the quadratic is tight \citep{ross2010efficient}. With
continuous actions it is worse:
\citet{simchowitz2025pitfalls} show that even under \eqref{eq:eiss} with
$\rho<1$ and a smooth deterministic expert, any smooth imitator can suffer
execution error $\exp(\Omega(T))$ times its error on the training distribution.
In the other direction \citet{zhang2025chunking} prove that chunking is the
remedy: once the chunk is longer than
$\log(1/\rho)^{-1}\log(\mathrm{poly}(L,C))$ the chunked closed loop is
exponentially stable again, with rate $\rho^{1/2}$, and its trajectory error
becomes horizon-free. That condition is a \emph{lower} bound on chunk length,
and it assumes a single $\rho<1$ for the system.

The bound in \eqref{eq:chunk} is the other side of the same inequality, and it
is why this paper measures rather than assumes. An admissible chunk length has to
satisfy both bounds at once,
\begin{equation}
\underbrace{\frac{\log \mathrm{poly}(L,C)}{\log(1/\rho)}}_{\text{fixed by the system}}
\;<\; k \;<\;
\underbrace{\frac{\log(1/\gamma)}{\lambda(s_t)}}_{\text{moves with the state}} ,
\label{eq:window}
\end{equation}
and Claim~1 is the finding that the right-hand side is not a constant: $\lambda$
takes both signs inside a single episode, so a single $k$ per task must respect
the tightest ceiling the episode ever imposes, leaving it conservative wherever
$\lambda$ is small and unsafe wherever $\lambda$ is large. The switching rule of
this paper is the two-level approximation of \eqref{eq:window}. Neither
\eqref{eq:chunk} nor \eqref{eq:window} is a theorem from the works cited; both
are the local reading of \eqref{eq:eiss} that a per-state measurement of
$\lambda$ makes available, and $\gamma$ in particular is a
modelling assumption about what one replan buys, which we do not measure.
Appendix~\ref{app:theory} gives the derivation, the statements quoted from, and
what the correspondence between $\lambda$ and $\log\rho$ does and does not
license.

\paragraph{Open-loop stability labels.}
For a state $s_{t_0}$ with actions $a_{t_0}, \dots, a_{t_0+K-1}$ we measure
divergence under action perturbation. First, restore the simulator to
$s_{t_0}$ and replay the $K$ actions unperturbed, recording a task metric
$y_j$, here end-effector position and gripper joint positions. Second, in $N$
branches, perturb the translation components of $a_{t_0}$ with noise of scale
$\sigma_u$, replay the same remaining actions, and record $y^{(n)}_j$. Third,
fit by least squares the slope $\lambda$ of
$\frac{1}{N}\sum_n \log \lVert y^{(n)}_j - y_j \rVert$ against $j$. This
$\lambda$ is a finite-time maximal Lyapunov exponent for the chunk. We use
$K{=}24$, $N{=}8$, a stride of 2, and $\sigma_u{=}0.05$ on RoboCasa. The label
is open-loop by construction: no replanning happens inside the window, so
$\lambda$ measures the plant and the task geometry under the recorded
behavior, not the reactivity of any controller. The procedure trains nothing
and reads no policy internals.
Figure~\ref{fig:labeling} shows the construction, the fit on two real states,
and the distribution the threshold below is read from.
Appendix~\ref{app:labeling} gives the procedure as pseudocode, the per-platform
task metric, and the reproducibility details, and Appendix~\ref{app:sigma}
covers how $\sigma_u$ is set and how much the labels depend on it.

% Labeling-algorithm figure: how one stability label is produced.
%
% Panel (a) is a schematic of the procedure in Sec. 3.
% Panels (b) and (c) are real data:
%   (b) mean-over-branches log divergence, regenerated for
%       rc_PickPlaceCounterToCabinet demo 20 by replaying the labeler's exact
%       call sequence and RNG stream (seed = demo index = 20) and keeping the
%       per-step curves that the production labeler discards.  The recomputed
%       slopes reproduce the published labels to <1e-4, so the lambda values
%       quoted here are the same numbers plotted in Fig. 1.
%   (c) histogram of lambda_task over all 23,724 stamps of that task, with the
%       automatic threshold delta = .0649 and the 5.0% of stamps above it.
% Regeneration script: scratchpad/dump_divergence.py

\begin{figure}[t]
\centering

%% ---------------------------------------------------------------- panel (a)
\begin{tikzpicture}[font=\footnotesize,
  stp/.style={circle, draw=cline, fill=white, inner sep=0.5pt,
              minimum size=2.7mm, font=\tiny, line width=0.4pt},
  sbox/.style={draw=cline, fill=cbox, rounded corners=1.6pt, align=center,
               inner sep=2.5pt, font=\tiny, line width=0.5pt},
  ar/.style={-latex, cline, line width=0.5pt},
]
% tau = 0 at x=2.4, tau = 24 at x=9.2 (0.2833 cm per step); nominal at y=1.75
\node[sbox, text width=10mm] (s0) at (1.05,1.75) {restore\\$s_{t_0}$};
\draw[ar] (s0.east) -- (2.36,1.75);
\node[stp] at (0.22,1.75) {1};

% the N=8 perturbed branches: same recorded actions, different first action
\draw[cline!60, line width=0.5pt] (2.4,1.75) .. controls (5.0,1.918) and (7.4,2.254) .. (9.2,2.87);
\draw[cline!60, line width=0.5pt] (2.4,1.75) .. controls (5.0,1.879) and (7.4,2.137) .. (9.2,2.61);
\draw[cline!60, line width=0.5pt] (2.4,1.75) .. controls (5.0,1.839) and (7.4,2.016) .. (9.2,2.34);
\draw[cline!60, line width=0.5pt] (2.4,1.75) .. controls (5.0,1.798) and (7.4,1.894) .. (9.2,2.07);
\draw[cline!60, line width=0.5pt] (2.4,1.75) .. controls (5.0,1.711) and (7.4,1.633) .. (9.2,1.49);
\draw[cline!60, line width=0.5pt] (2.4,1.75) .. controls (5.0,1.671) and (7.4,1.511) .. (9.2,1.22);
\draw[cline!60, line width=0.5pt] (2.4,1.75) .. controls (5.0,1.628) and (7.4,1.386) .. (9.2,0.94);
\draw[cline!60, line width=0.5pt] (2.4,1.75) .. controls (5.0,1.578) and (7.4,1.233) .. (9.2,0.60);

% the nominal branch
\draw[cline, line width=0.9pt] (2.4,1.75) -- (9.2,1.75);
\fill[cline] (2.4,1.75) circle (1.3pt);

% what makes the branches differ
\node[font=\tiny, cunstable, anchor=south west, align=left, inner sep=1pt]
  (pert) at (2.95,2.45)
  {perturb $a_{t_0}$ only: $+\,\sigma_u\varepsilon$\\
   on the translation dims};
\draw[-latex, cunstable, line width=0.5pt] (3.05,2.43) -- (2.63,1.83);
\node[font=\tiny, cline, anchor=south west, inner sep=1pt]
  at (2.60,0.62) {identical actions in every branch after the first};

% the quantity that is measured
\draw[latex-latex, cline, line width=0.6pt] (7.5,1.75) -- (7.5,2.40);
\node[font=\tiny, fill=white, inner sep=1pt] at (7.5,2.08) {$d_n(\tau)$};
\node[stp] at (7.5,2.62) {4};

\node[stp] at (9.38,1.75) {2};
\node[font=\tiny, anchor=west, inner sep=1.5pt] at (9.52,1.75) {nominal};
\node[stp] at (9.38,2.87) {3};
\node[font=\tiny, anchor=west, inner sep=1.5pt, align=left]
  at (9.52,2.87) {$N{=}8$\\branches};

% tau axis
\draw[-latex, cline, line width=0.5pt] (2.4,0.30) -- (9.55,0.30);
\draw[cline, line width=0.5pt] (2.4,0.30) -- (2.4,0.21);
\draw[cline, line width=0.5pt] (4.1,0.30) -- (4.1,0.21);
\draw[cline, line width=0.5pt] (5.8,0.30) -- (5.8,0.21);
\draw[cline, line width=0.5pt] (7.5,0.30) -- (7.5,0.21);
\draw[cline, line width=0.5pt] (9.2,0.30) -- (9.2,0.21);
\node[font=\tiny, anchor=north, inner sep=2pt] at (2.4,0.21) {0};
\node[font=\tiny, anchor=north, inner sep=2pt] at (4.1,0.21) {6};
\node[font=\tiny, anchor=north, inner sep=2pt] at (5.8,0.21) {12};
\node[font=\tiny, anchor=north, inner sep=2pt] at (7.5,0.21) {18};
\node[font=\tiny, anchor=north, inner sep=2pt] at (9.2,0.21) {24};
\node[font=\tiny, anchor=west, inner sep=2pt] at (9.55,0.22) {$\tau$};

\node[anchor=north west, font=\footnotesize\bfseries] at (0.0,3.10) {(a)};
% the box is the full text width, and the (b)+(c) row below is stretched to
% the same width with \hfill, so the two rows share a left edge
\useasboundingbox (0.0,-0.08) rectangle (\textwidth,3.10);
\end{tikzpicture}

\vspace{2pt}

%% ---------------------------------------------------------------- panel (b)
\begin{tikzpicture}[font=\footnotesize]
\begin{axis}[
  scale only axis, width=4.5cm, height=2.6cm,
  xmin=1, xmax=24, ymin=-7.2, ymax=-3.0,
  xtick={1,6,12,18,24}, ytick={-7,-6,-5,-4},
  ticklabel style={font=\scriptsize},
  xlabel={$\tau$}, ylabel={mean $\log d_n(\tau)$},
  xlabel style={font=\scriptsize, yshift=3pt},
  ylabel style={font=\scriptsize, yshift=-4pt},
  axis line style={cline}, tick style={cline}, tick align=outside,
  every axis plot/.append style={line width=0.9pt},
  clip mode=individual,
]
  % t0=48: unstable, d grows 0.0018 m -> 0.0256 m over the window
  \addplot[cunstable] coordinates {
    (1,-6.3076) (2,-5.9927) (3,-6.2816) (4,-5.5934) (5,-5.3721) (6,-5.2209)
    (7,-5.1520) (8,-5.1344) (9,-5.1184) (10,-5.0567) (11,-5.1292) (12,-4.8690)
    (13,-4.5927) (14,-4.3482) (15,-4.1399) (16,-3.9602) (17,-3.8317)
    (18,-3.7741) (19,-3.7472) (20,-3.7123) (21,-3.6822) (22,-3.6773)
    (23,-3.6709) (24,-3.6637)
  };
  \addplot[cunstable, dashed, line width=0.6pt]
    coordinates {(1,-6.0568) (24,-3.2789)};
  % t0=10: stable, d stays at 0.0013-0.0018 m
  \addplot[cstable] coordinates {
    (1,-6.6133) (2,-6.4940) (3,-6.4810) (4,-6.4753) (5,-6.3353) (6,-6.2190)
    (7,-6.1582) (8,-6.1909) (9,-6.2181) (10,-6.2024) (11,-6.1560) (12,-6.1454)
    (13,-6.0897) (14,-6.0891) (15,-6.1115) (16,-6.1230) (17,-6.1617)
    (18,-6.2319) (19,-6.2902) (20,-6.3413) (21,-6.3524) (22,-6.3475)
    (23,-6.3614) (24,-6.3434)
  };
  \addplot[cstable, dashed, line width=0.6pt]
    coordinates {(1,-6.3393) (24,-6.2050)};
  \node[anchor=north west, font=\tiny, cunstable, align=left, inner sep=1pt]
    at (axis cs:1.4,-3.02) {unstable state, $\lambda{=}.121$};
  \node[anchor=south west, font=\tiny, cstable, align=left, inner sep=1pt]
    at (axis cs:1.4,-7.18) {stable state, $\lambda{=}.006$};
\end{axis}
\node[anchor=north west, font=\footnotesize\bfseries] at (-1.15,2.95) {(b)};
\end{tikzpicture}
\hfill
%% ---------------------------------------------------------------- panel (c)
\begin{tikzpicture}[font=\footnotesize]
\begin{axis}[
  scale only axis, width=4.5cm, height=2.6cm,
  xmin=-0.10, xmax=0.16, ymin=0, ymax=19.6,
  xtick={-0.1,0,0.1}, ytick={0,5,10,15},
  ticklabel style={font=\scriptsize},
  xlabel={$\lambda$ over all stamps of the task},
  ylabel={\% of stamps},
  xlabel style={font=\scriptsize, yshift=3pt},
  ylabel style={font=\scriptsize, yshift=-4pt},
  axis line style={cline}, tick style={cline}, tick align=outside,
  clip mode=individual,
]
  \fill[cband] (axis cs:0.0649,0) rectangle (axis cs:0.16,19.6);
  \addplot[ybar interval, draw=cstable, fill=cstable!25, line width=0.3pt]
    coordinates {
    (-0.10,0.510) (-0.09,0.691) (-0.08,0.784) (-0.07,0.906) (-0.06,1.290)
    (-0.05,1.960) (-0.04,3.267) (-0.03,6.171) (-0.02,10.538) (-0.01,16.051)
    (0.00,18.230) (0.01,13.674) (0.02,7.478) (0.03,5.105) (0.04,3.971)
    (0.05,3.216) (0.06,2.234) (0.07,1.467) (0.08,1.113) (0.09,0.577)
    (0.10,0.375) (0.11,0.181) (0.12,0.122) (0.13,0.046) (0.14,0.025)
    (0.15,0.013) (0.16,0)
  };
  \draw[cunstable, dashed, line width=0.7pt]
    (axis cs:0.0649,0) -- (axis cs:0.0649,19.6);
  \node[anchor=north east, font=\tiny, cunstable, inner sep=1pt]
    at (axis cs:0.0625,19.4) {$\delta$};
  \node[anchor=north east, font=\tiny, cunstable, align=right, inner sep=1pt]
    at (axis cs:0.158,14.0) {unstable:\\5.0\% of\\stamps};
\end{axis}
\node[anchor=north west, font=\footnotesize\bfseries] at (-1.15,2.95) {(c)};
\end{tikzpicture}

\caption{\textbf{How one stability label is produced.}
\textbf{(a)} The procedure at a single state, and the whole of it: restore the
simulator to $s_{t_0}$ (1); replay the $K{=}24$ recorded actions unperturbed to
get the nominal branch (2); repeat $N{=}8$ times with Gaussian noise added to
the translation components of the first action only, every branch replaying the
same recorded actions thereafter (3); and record the task-space distance
$d_n(\tau)$ between each branch and the nominal at every step (4). The label
$\lambda$ is the least-squares slope of the mean log distance against $\tau$.
Nothing is trained, no policy internals are read, and the policy is not involved
at all: the label describes the plant and the task geometry under the recorded
behavior.
\textbf{(b)} The fit, on two real states drawn from the same episode as
Figure~\ref{fig:framework}, with the least-squares line dashed. At the unstable
state the branches separate from 1.8 to 25.6 millimetres in the task metric
over the window, a factor of 14; at the stable state they stay between 1.3 and
1.8 millimetres,
and the fitted slope is twenty times smaller. The label is the slope,
not the final distance, so it does not reward states that merely start far
apart. Neither curve is a clean straight line, and the fit is a summary of the
window rather than a claim that growth is exactly exponential within it.
\textbf{(c)} Where the threshold comes from. Distribution of $\lambda$ over all
23{,}724 stamps of PickPlaceCounterToCabinet. The bulk sits at a task-specific
noise floor set by contact chatter and metric scale, so a fixed cutoff would not
transfer between tasks; $\delta$ is instead read off the below-median half of
this distribution as $P_{95}(\lvert\lambda\rvert)$, which puts 5.0\% of stamps
in the unstable class here. Stamps with $\lvert\lambda\rvert \le \delta$ are the
deadband and are excluded from classification training.}
\label{fig:labeling}
\end{figure}


\paragraph{Automatic threshold.}
Divergence slopes have a task-specific noise floor from contact chatter and
metric scale. We set
$\delta = P_{95}\left(\lvert \lambda \rvert \,:\, \lambda < \mathrm{median}(\lambda)\right)$,
the 95th percentile of the magnitude of the below-median half, and call a stamp
unstable if $\lambda > \delta$. This adapts to each task's floor with no manual
tuning. Stamps with $\lvert \lambda \rvert \le \delta$ form a deadband where
the label is ambiguous by construction, and we exclude them from
classification training and report classification quality both with and
without them.

\paragraph{Closed-loop stability labels.}
The same construction gives a second axis if the perturbed branch is driven by
the policy replanning at every step instead of replaying fixed actions. Two
choices matter. The nominal branch must also be policy-driven and unperturbed,
otherwise the measured rate mixes perturbation growth with the ordinary
difference between the policy and the recorded behavior. And because modern
policies sample stochastically, two unperturbed policy rollouts from one state
already separate, so we run a control pair from the same state with no
perturbation and report its rate $\lambda_{\text{ctrl}}$ alongside. A segment
is closed-loop stable when the perturbation adds nothing beyond that floor.

\paragraph{Two axes, four regimes.}
Open-loop stability says whether errors grow while a chunk plays out.
Closed-loop stability says whether the policy tolerates frequent replanning,
since some tasks degrade sharply at $k{=}1$ when replanning every step
introduces dithering near fine goals. The two axes give four regimes.
Open-loop stable and closed-loop stable: any $k$ works, prefer long chunks.
Open-loop unstable and closed-loop stable: the case that motivates switching,
because short chunks fix what long chunks break.
Open-loop stable and closed-loop unstable: long chunks are mandatory and a
switcher must not produce false positives here.
Unstable on both axes: no chunk length helps and the failure is not a chunking
problem.

\paragraph{Predictors.}
To act on the signal at rollout time without a simulator we train a small
attentive head on frozen V-JEPA features that regresses $\lambda$ and
classifies $\lambda > \delta$. The h1 variant reads one clip of 16 consecutive
frames, about 0.8 seconds, together with the robot state over the same window.
The h4 variant reads four such clips spaced 16 frames apart, about 3.2 seconds,
with a learned per-clip position embedding, and tests whether the regime is
visible in motion over time. We also evaluate a simulator oracle that runs the
labeling procedure live at each replan and therefore measures $\lambda$ instead
of predicting it. The oracle needs to roll the simulator forward and rewind it,
so it cannot run on a real robot; its role is to bound what perfect prediction
would be worth. Appendix~\ref{app:predictor} gives the architecture and
training setup.

\paragraph{Switching rule.}
At each replan the controller obtains a stability score and executes
$k = k_u$ if the score exceeds a threshold, else $k = k_s$, with
$(k_s, k_u) = (16, 1)$ throughout unless stated. The rule only shortens on
predicted instability, so in stable segments it costs nothing by construction.
Setting the threshold for a learned head is not trivial and we treat it as part
of the method rather than an implementation detail; Section~\ref{sec:claim3}
reports the two ways it fails and the calibration we adopt.
